An der Hochschule sind Sitzmöglichkeiten, Platz zum Schreiben und auch Schreibmaterialien sowie Toiletten, Klimatisierung und eine ausreichende Beleuchtung vorhanden. Der Zugang zu grundlegender elektronischer Infrastruktur ist gewährleistet, dazu zählen insbesondere Steckdosen und eine zeitgemäße Internetverbindung. Diese Infrastruktur ist während der Vorlesungszeiten für die Studierenden barrierefrei zugänglich.

Aufeinander folgende Pflichtmodule sind räumlich nahe beieinander, damit genug Zeit zwischen den Veranstaltungen besteht, um an den nächsten Veranstaltungsort zu gelangen. Dabei sind die Gebäude, in denen die Veranstaltung stattfinden, deutlich sichtbar und auch international verständlich gekennzeichnet. Zur Orientierung befinden sich an zentralen Stellen Übersichtspläne.

Es gibt eine Druckmöglichkeit höchstens zum Selbstkostenpreis. Die Funktionsfähigkeit dieser ist gewährleistet (Toner, Papier vorhanden).

\section{Räume}

Die Räumlichkeiten der Hochschule sind barrierefrei erreichbar. Je nach Verwendungszweck können sich die Anforderungen an die Räume unterscheiden, daher wird hier in drei für Studierende relevante Raumarten unterschieden.

\subsection{Veranstaltungsräume}

Für jede Veranstaltung steht ein Raum zur Verfügung, der über hinreichend viele Sitzplätze verfügt. Auch zu Stoßzeiten sind ausreichend Kapazitäten an Räumlichkeiten vorhanden. Es gibt Platz, um Jacken, Taschen usw. abzulegen.

Die Räume verfügen über eine Tafel bzw. ein Whiteboard in hinreichender Größe, sodass bei einer für alle Anwesenden lesbar großen Anschrift genügend Tafelfläche vorhanden ist, um den für das Verständnis des aktuellen Themas notwendigen Kontext zu fassen. Außerdem bieten die Räumlichkeiten in der Regel die Möglichkeit zur Visualisierung per Beamer. Dafür sind Projektionsflächen vorhanden. In den Räumen sind die Dozierenden überall zu verstehen; geeignete Hilfsmittel, wie bspw. Mikrofone, stehen bei Bedarf zur Verfügung.

\subsection{Computerräume}

Die Anzahl frei zugänglicher Rechnerplätze ist der Anzahl der Studierenden angemessen. Für Veranstaltungen mit PC-Einsatz gibt es genügend PCs, sodass sich maximal zwei Studierende einen PC teilen müssen.

Die Rechner sind entsprechend der Richtung der Hochschule mit Software ausgestattet (Computeralgebrasystem, Numerische Software, Statistikprogramm, \dots). Ein Programm zum Anfertigen auch umfangreicherer mathematischer Texte ist installiert. Die Rechner verfügen über einen Internetzugang.

Es gibt eine Möglichkeit zur Visualisierung bei Lehrveranstaltungen (z.B. Beamer, Tafel, \dots).

\subsection{Studentische Arbeitsräume}

Für die Studierende besteht die Möglichkeit sich über die Belegung von Räumen zu informieren und freie Räume als Arbeitsraum zu nutzen. Studierende, die an ihrer Abschlussarbeit arbeiten, haben immer Zugang zu einem Rechner. Studentische Arbeitsräume sind mit einer Tafel bzw. einem Whiteboard ausgestattet.

\section{Fachschaft}

Die Fachschaft ist ein Zusammenschluss von Studierenden eines Faches oder mehrerer verwandter Fächer zwecks gemeinsamer Interessenvertretung. Die Fachschaft wird durch die Hochschule unterstützt.

Um eine effiziente Fachschaftsarbeit zu gewährleisten, stellt die Hochschule der Fachschaft einen Büroraum, dessen Größe der Anzahl der Studierenden angemessen und gut erreichbar ist, zur Verfügung. Ebenso werden ein Telefon, ein Rechnerzugang inklusive Webspace für eine Fachschaftswebseite und eine E-Mail-Adresse sowie eine Kopier- und Druckmöglichkeit bereitgestellt.

Außerdem ermöglicht die Hochschule der Fachschaft regelmäßige Fachschaftssitzungen und die Herausgabe von Infomaterial, z.B. den Druck eines regelmäßigen Infohefts, von Plakaten, Flyern und ähnlichem.

\section{Bibliothek}

Eine Bibliothek ist vorhanden und verfügt über den Veranstaltungen zugrunde liegende und vertiefende Literatur. Dazu gehören Bücher und (Fach-)Zeitschriften in analoger und auch digitaler Form. Sind diese weder analog noch digital verfügbar, kann die Literatur über die Fernleihe von anderen Bibliotheken bezogen werden. Der Bestand ist digital durchsuchbar und wird regelmäßig im Rahmen der Bedarfüberprüfung aktualisiert.

An der Bibliothek sind auch Arbeitsplätze verfügbar, ebenso ein Kopierer bzw. Scanner und Recherchemöglichkeiten.

\section{Digitale Infrastruktur}

(Immatrikulations-)Bescheinigungen sind online ausstellbar und studiumsrelevante Ordnungen und Formulare sind digital auffindbar. Die Prüfungsanmeldung ist online und barrierefrei durchführbar. Ebenso gibt es eine digitale Lernplattform für Vorlesungen, weitere Informationen, etc. und ein Mail-Postfach mit sinnvoll bemessenen Speicherplatz.

Für das Studium notwendige Software wird von der Hochschule bereitgestellt. Neben fachspezifischer Software gehört dazu ein Cloudspeicherzugang, eine Videokomunikationsplattform sowie eine Plattform für kollaboratives Arbeiten.

Der Zugang zu geeigneten Rechnern, insbesondere für Abschlussarbeiten, wird ermöglicht, bspw. durch die Möglichkeit Laptops auszuleihen. Es gibt eine Ansprechstelle für IT-Probleme.

Auf eine sichere und datenschutzkonforme Umsetzung wird geachtet.
